% Generated draft: requires mathematical and source review.
\documentclass[11pt]{article}
\usepackage[margin=25mm]{geometry}
\usepackage{fontspec}
\setmainfont{Latin Modern Roman}
\usepackage{amsmath,amssymb,mathtools,unicode-math}
\setmathfont{Latin Modern Math}
\usepackage{enumitem,needspace,xurl,hyperref}
\setlength{\parindent}{0pt}
\setlength{\parskip}{5pt}
\setlength{\emergencystretch}{3em}
\newcommand{\N}{\mathbb{N}}
\newcommand{\Z}{\mathbb{Z}}
\newcommand{\Q}{\mathbb{Q}}
\newcommand{\R}{\mathbb{R}}
\newcommand{\C}{\mathbb{C}}
\newcommand{\F}{\mathbb{F}}
\newcommand{\A}{\mathbb{A}}
\newcommand{\PP}{\mathbb{P}}
\newcommand{\E}{\mathbb{E}}
\newcommand{\eps}{\varepsilon}
\newcommand{\dd}{\,\mathrm{d}}
\begin{document}
\section*{UnsolvedMath ID: 30006257}
\textbf{Source problem code:} catalogue entry
\subsection{Definitions and mathematical statement}
The supplied record concerns a small quantum group and the blocks of its center. Write \(u_\lambda\) for the algebra block associated with a weight \(\lambda\), and write \(z_\lambda=Z(u_\lambda)\) for its center. Given the intended action of an algebraic group \(G\) on this center, define
\[
z_\lambda^G=\{c\in z_\lambda: g\cdot c=c\text{ for every }g\in G\}.
\]
Let \(W\) be the Weyl group in the intended setting and let
\[
W_P=\operatorname{Stab}_W(\lambda).
\]
The action used to form this stabilizer is not specified in the supplied record.

For clarity, the usual diagonal coinvariant algebra associated with a reflection representation \(\mathfrak h\) over \(\C\) is
\[
D_W=\C[\mathfrak h\oplus\mathfrak h^*]\big/\big\langle\C[\mathfrak h\oplus\mathfrak h^*]^W_+\big\rangle,
\]
where \(W\) acts diagonally, the subscript \(+\) denotes invariant polynomials with zero constant term, and the brackets denote the ideal they generate. Its polynomial bigrading records degrees in the two summands separately. The source calls \(DR_W^{W_P}\) a canonical quotient of the diagonal coinvariant algebra and says that it is a module over a Cherednik algebra. However, the displayed defining expression in the supplied original statement is corrupted, particularly the placement of \(W_P\). Consequently, the precise quotient denoted by \(DR_W^{W_P}\) cannot be recovered unambiguously from this record; the notation here retains the source's intended object without identifying it with \(D_W\) or with an invariant subspace of \(D_W\).

\textbf{Conjectural target, conditional on the missing specifications.} For every weight \(\lambda\) covered by the source's small quantum group setting and its restrictions on \(l\), is there an isomorphism
\[
z_\lambda^G\cong DR_W^{W_P}
\]
as bigraded complex vector spaces? Such an isomorphism must preserve each bidegree. The record does not specify the admissible quantum group parameters, the restrictions on \(l\), the allowed weights, the group action, the bigrading on the left, or any grading conversion between the two sides. Thus this is a conditional formulation of the intended question, not a fully specified conjecture for arbitrary small quantum groups.

The surrounding excerpt also contains a purported commutative algebra identification of \(z_\lambda\) with equivariant Hochschild cohomology and a sum of polyvector-field cohomology groups subject to \(i+j+k=0\). Its variety notation and displayed isomorphisms are corrupted, and its hypotheses are absent. It therefore cannot supply an additional precise assertion or a definition of the required bigrading here.
\subsection{Short English statement}
The question asks whether the \(G\)-fixed part of each relevant small quantum group center block has the same bigraded vector-space structure as the source's parabolic diagonal coinvariant quotient. The quotient, grading conventions, and restrictions on the quantum group parameter must be recovered before the conjecture is fully precise. The supplied status is partially solved; this draft does not independently verify that assessment.
\subsection{Sources}
\begin{itemize}
\item Local registry: \texttt{data/unsolvedmath/problems.json}, with \texttt{catalog/problems.json} as the additional-record registry; permanent UnsolvedMath ID \(30006257\), title \emph{Centers of Small Quantum Groups}, exact source code \texttt{catalogue entry}. The supplied complete record is the evidence used for this draft.
\item Supplied authoritative report citation: \emph{Algebraic Groups} (2025), Oberwolfach Reports, pp. 941--1010, DOI \texttt{10.4171/owr/2025/20}. \url{https://doi.org/10.4171/owr/2025/20}. The excerpt labels the target Conjecture 3 and attributes its formulation to studies cited there as [6, 7], with some restrictions on \(l\); it cites [4] for the Cherednik algebra connection. The corresponding numbered bibliography entries are not supplied. No individual proposer is named in the record.
\item Supplied literature source, OWR 2025 report: \url{https://ems.press/content/serial-article-files/51849}.
\item Supplied literature source, Lachowska--You, principal blocks: \url{https://doi.org/10.1093/imrn/rnx062}.
\item Supplied literature source, Lachowska--You, singular blocks: \url{https://doi.org/10.1112/plms.12193}.
\item Literature assessment dated exactly \texttt{2026-08-22}: the supplied curation labels the problem partially solved and reports a dimension equality and substantial type-A results, while stating that the bigraded isomorphism is not proved in general. This is supplied dated curation, not an independently verified literature assessment. The record itself requests independent verification of time-sensitive status and cited claims.
\item The registry labels its statement \texttt{corrected\_verified} and reports that the official report was checked and its notation repaired. These are supplied source-verification claims; no source visit or independent verification was performed for this draft.
\end{itemize}
\end{document}
